\introChapter{ทฤษฎีเซต (Set Theory)}

\begin{description}[leftmargin=2cm, style=nextline, itemsep=0.1em, parsep=0pt]
    \item[รหัสวิชา] 4091606
    \item[ชื่อวิชา] คณิตศาสตร์สำหรับคอมพิวเตอร์
    \item[หน่วยกิต] 3 (3-0-6)
    \item[บทที่] 2
    \item[ชื่อบท] ทฤษฎีเซต (Set Theory)
    \item[จำนวนชั่วโมง] 6 ชั่วโมง
\end{description}

\noindent \textbf{\large วัตถุประสงค์การเรียนรู้ประจำบท}\par\vspace{0.2em}\noindent
\noindent เมื่อนักศึกษาเรียนบทนี้จบแล้ว นักศึกษาจะสามารถ
\begin{enumerate}[topsep=0.2em, itemsep=0.15em, leftmargin=*]
    \item อธิบายความหมายและสัญลักษณ์ของเซตได้
    \item ระบุประเภทของเซต เช่น เซตว่าง เซตจำกัด เซตอนันต์ได้
    \item ใช้การดำเนินการบนเซต เช่น ยูเนียน อินเตอร์เซกชัน ผลต่างเซต และคอมพลีเมนต์ได้
    \item พิสูจน์สมบัติของเซตโดยใช้แผนภาพเวนน์และพีชคณิตเซตได้
    \item ประยุกต์ใช้ทฤษฎีเซตในการแก้ปัญหาทางคอมพิวเตอร์ได้
\end{enumerate}

\noindent \textbf{\large หัวข้อที่สอน}\par\vspace{0.2em}\noindent
\begin{enumerate}[topsep=0.2em, itemsep=0.15em, leftmargin=*]
    \item นิยามและสัญลักษณ์ของเซต
    \item ความสัมพันธ์ระหว่างเซตและประเภทของเซต
    \item การดำเนินการระหว่างเซตและแผนภาพเวนน์
    \item กฎของเดมอร์แกน
    \item การนับจำนวนสมาชิกของเซต
    \item เซตที่นับได้และเซตที่นับไม่ได้
    \item การประยุกต์ทฤษฎีเซตในวิทยาการคอมพิวเตอร์
\end{enumerate}

\noindent \textbf{\large สื่อการสอน}
\begin{enumerate}[topsep=0.2em, itemsep=0.15em, leftmargin=*]
    \item สไลด์ประกอบการสอน
    \item แผนภาพเวนน์
    \item แบบฝึกหัดท้ายบท
\end{enumerate}
